\subsubsection{Đánh giá kiểm định các cặp mô hình}

\tablename~\ref{tab:classification-mcnemar-test} thể hiện được kết quả kiểm
định McNemar các cặp mô hình về sự giống nhau của phân phối lỗi. Các kết quả
được làm tròn đến 4 chữ số thập phân.

\begin{table}[htbp]
  \centering
  \caption{Kết quả kiểm định McNemar các cặp mô hình trên tập kiểm tra}
  \label{tab:classification-mcnemar-test}
  \import{\tablespath}{classification-mcnemar-test}
\end{table}

Kết quả kiểm định McNemar trong
\tablename~\ref{tab:classification-mcnemar-test} cho thấy tất cả các cặp mô
hình đều có p-value xấp xỉ  0.0000 (sau khi làm tròn 4 chữ số thập phân). Điều
này cho phép bác bỏ giả thuyết $H_0$ ($H_0$ là  hai mô hình có cùng phân phối
lỗi) ở mức ý nghĩa thông thường (ví dụ 0.05). Nói cách khác, sự khác biệt về
hiệu năng giữa mọi cặp mô hình đều có ý nghĩa thống kê, không phải do ngẫu
nhiên.

Đáng chú ý, các cặp liên quan đến Softmax có giá trị thống kê rất lớn (đặc biệt
với Weighted Softmax: 8905.0), cho thấy sự khác biệt rõ rệt trong cách hai mô
hình này dự đoán. Điều này phù hợp với kết quả trước đó: dù cùng họ mô hình,
Softmax tối ưu Accuracy trong khi Weighted Softmax thiên về Recall, dẫn đến sự
khác biệt đáng kể trong các dự đoán sai/đúng.

Tương tự, các cặp như LDA – QDA (5508.0) hay LDA – Weighted Softmax (6236.0)
cũng có thống kê lớn, phản ánh sự khác biệt rõ ràng về bản chất mô hình và hành
vi phân loại. Đặc biệt, dù QDA có Recall tương đối cao, kết quả kiểm định cho
thấy mô hình này vẫn tạo ra mẫu lỗi rất khác biệt so với các mô hình còn lại,
phù hợp với việc Accuracy và F1 của QDA thấp.

Tổng thể, kiểm định McNemar củng cố kết luận từ các bảng trước: các mô hình
không chỉ khác nhau về giá trị chỉ số mà còn khác nhau một cách có ý nghĩa
thống kê trong dự đoán thực tế. Điều này khẳng định rằng việc lựa chọn mô hình
(Softmax, LDA, hay Weighted Softmax) sẽ dẫn đến những hành vi phân loại khác
biệt rõ rệt, tùy thuộc vào mục tiêu tối ưu (Accuracy, Recall hay cân bằng tổng
thể).
